% $Id: vortrag.tex,v 1.1 2001/05/06 11:43:21 thor Exp $
%
\documentclass{slides}
\usepackage{times}
%\usepackage[light]{draftcopy}
%\usepackage{a4}
\usepackage{german}
\usepackage[latin1]{inputenc}
%\usepackage{amsmath}
%\usepackage{amssymb}
\usepackage{xspace}
\usepackage{graphicx}
\usepackage{fancyhdr}
\usepackage{xspace}
\usepackage{verbatim}
\usepackage{url}
%\usepackage[notoday,nofancy]{rcsinfo}
%\rcsInfo $Id: vortrag.tex,v 1.1 2001/05/06 11:43:21 thor Exp $
%\parindent0ex
%\headheight15pt
%\selectlanguage{german}

\newcommand{\ZIB}{\textsc{Zib}\xspace}
\newcommand{\CPlex}{\textsc{CPlex}\xspace}
\newcommand{\Zimpl}{\textsc{Zimpl}\xspace}

\renewcommand{\footrulewidth}{\headrulewidth}
\rhead{\includegraphics[scale=.4,bb=200 110 230 190]{ziblogo2.eps}}
%\epsfig{file=ziblogo2.eps,width=2cm}}
\lfoot{\tiny Mathematical Modeling Languages}
\cfoot{}
\pagestyle{fancy}
\headheight60pt
\headsep10pt
\begin{document}
\begin{slide}
\lhead{Mathematical Modeling Languages}
\thispagestyle{fancy}

What is a MPL?

\small
Mathematical or algebraic modeling languages allow to formulate a
linear or nonlinear programming problem in 
a way near the original mathematical formulation and 
automatically generate data for a corresponding solver.

\normalsize Uses for MPL's

\small
\begin{itemize}
\item Teaching modeling
\item Rapid prototyping
\item Projects with industry
\end{itemize}

\normalsize Some links

\small
NEOS Guide\\
\footnotesize
\url{http://www-fp.mcs.anl.gov/otc/Guide/}\\
\url{SoftwareGuide/Categories/optsysmodlang.html}

\small
Arnold Neumaier\\
\footnotesize
\url{http://solon.cma.univie.ac.at/~neum/glopt.html}

\small
Decision Tree for Optimization Software\\
\footnotesize
\url{http://plato.la.asu.edu/topics/tools.html}

\vfill
\end{slide}

\begin{slide}
\lhead{Some MPL's}
\thispagestyle{fancy}

\vfill
\tiny

\begin{tabular}{llll}
%------------------------------------------------------------------------------
Name&Author&Models\\
\hline\\
%------------------------------------------------------------------------------
AIMMS&
Johannes J. Bisschop&
LP, MIP, NLP, NWP\\
&Roobert Entriken\\
&Paragon Decision\\
&Technology\\
\\\hline\\
%------------------------------------------------------------------------------
AMPL&
Robert Fourer&
LP, MIP, QP, NLP, NWP\\
&David M. Gay\\
&Brian W. Kernighan\\
&AT\&T Bell Laboratories\\
\\\hline\\
%------------------------------------------------------------------------------
GAMS&
Anthony Brooke&
LP, MIP, QP, NLP, NWP\\
&David Kendrick\\
&Alexander Meeraus\\
&Ramesh Raman\\
&GAMS Development\\
&Corporation\\
\\\hline\\
%------------------------------------------------------------------------------
LINGO&
LINDO Systems&
LP, MIP, QP, NLP\\
\\\hline\\
%------------------------------------------------------------------------------
MINOPT&
C. A. Schweiger&
LP, MIP, NLP, MINLP\\
&C. A. Floudas\\
&Princeton University\\
\\\hline\\
%------------------------------------------------------------------------------
OPL&
Pascal van Hentenryck&
LP, MIP, NLP\\
&Irvin Lustig\\
&Larent Michael\\
&Jean-Fran\c{c}is Puget\\
&ILOG\\
\\\hline\\
%------------------------------------------------------------------------------
XPRESS-MP&
Dash Associates&
LP, MIP, NLP\\

\end{tabular}
\end{slide}

\begin{slide}
\lhead{AMPL Model}
\thispagestyle{fancy}
\tiny
\begin{verbatim}
set P;                # Products
set R;                # Raw materials

param T > 0;          # number of production periods
param M > 0;          # Maximum production per period
param a{R,P} >= 0;    # units of raw material i to 
                      # manufacture 1 unit of product j
param b{R} >= 0;      # maximum initial stock 
param c{P,1..T};      # estimated profit per unit 
param d{R};           # storage cost per period per unit 
param f{R};           # estimated remaining value per 
                      # unit after last period
var x{P,1..T} >= 0;   # units of product manufactured 
var s{R,1..T+1} >= 0; # stock of raw material 

maximize profit:
    sum {t in 1..T} (sum {j in P} c[j,t]*x[j,t] 
                 - sum {i in R} d[i]*s[i,t])
  + sum {i in R} f[i] * s[i,T+1];

subject to 

   prd {t in 1..T}: sum {j in P} x[j,t] <= M;
 
   stock1 {i in R}: s[i,1] <= b[i];

   stock {i in R, t in 1..T}: 
      s[i,t+1] = s[i,t] - sum {j in P} a[i,j] * x[j,t];
\end{verbatim}
\vfill
\end{slide}

\begin{slide}
\lhead{AMPL Data}
\thispagestyle{fancy}
\tiny
\begin{verbatim}

data;

set P := lite bud mich ;
set R := malt hops ;

param T 3 ;
param M 40 ;
param a [malt,lite] 5 [malt,bud] 3 [malt,mich] 1
        [hops,lite] 1 [hops,bud] 2 [hops,mich] 3 ;
param b [malt] 400 [hops] 275 ;
param c [lite,1] 25 [lite,2] 20 [lite,3] 10
        [bud,1] 50 [bud,2] 50 [bud,3] 50
        [mich,1] 75 [mich,2] 80 [mich,3] 100 ;
param d [malt] 0.5 [hops] 2.0 ;
param f [malt] 15 [hops] 25 ;
\end{verbatim}
\vfill
\end{slide}


\begin{slide}
\lhead{GAMS}
\thispagestyle{fancy}
\tiny
\begin{verbatim}
SETS
     I   canning plants / SEATTLE, SAN-DIEGO /
     J   markets        / NEW-YORK, CHICAGO, TOPEKA / ;
PARAMETERS
     A(I)  capacity of plant i in cases
       /    SEATTLE     350
            SAN-DIEGO   600  /
     B(J)  demand at market j in cases
       /    NEW-YORK    325
            CHICAGO     300
            TOPEKA      275  / ;
TABLE D(I,J)  distance in thousands of miles
                  NEW-YORK       CHICAGO      TOPEKA
    SEATTLE          2.5           1.7          1.8
    SAN-DIEGO        2.5           1.8          1.4  ;
SCALAR F  freight in USD per case per 1000 miles  /90/ ;
PARAMETER C(I,J) cost in 1000s of USD per case;
          C(I,J) = F * D(I,J) / 1000 ;
VARIABLES
     X(I,J)  shipment quantities in cases
     Z       total transportation costs in 1000s of USD;
POSITIVE VARIABLE X ;
EQUATIONS
     COST        define objective function
     SUPPLY(I)   observe supply limit at plant i
     DEMAND(J)   satisfy demand at market j ;
COST ..        Z  =E=  SUM((I,J), C(I,J)*X(I,J)) ;
SUPPLY(I) ..   SUM(J, X(I,J))  =L=  A(I) ;
DEMAND(J) ..   SUM(I, X(I,J))  =G=  B(J) ;
MODEL TRANSPORT /ALL/ ;
SOLVE TRANSPORT USING LP MINIMIZING Z ;
\end{verbatim}
\vfill
\end{slide}

\begin{slide}
\lhead{MINOPT}
\thispagestyle{fancy}
\tiny
\begin{verbatim}
DECLARATION {{
INDEX {i,j,k,p};
SET I = |1,2|;
SET J = |1,2|;
SET K = |1,2|;

PARAM c(I,J,K,P) = {10,15,20,10,5,5,30,10,25,20,15,20};
PARAM s(I,P) = {100,200,150,120};
PARAM f(J) = {2000,2500};
PARAM v(J) = {50,70};

XVAR {x(I,J,K,P)};
YVAR {y(J,K),z(J)};
BINARY {y(J,K),z(J)};
POSITIVE {x(I,J,K,P)};
UBDS x(I,J,K,P)=<i E I|<j E J|<k E K|<p E P| 300>>>>;
YSTP {1,0,0,1,1,1};
}}
MODEL {{
  MIN: <<i E I| <<j E J| <<k E K| 
          <<p E P| c(i,j,k,p)*x(i,j,k,p) >> >> >> >>
       + <<j E J| f(j)*z(j) >>
       + <<j E J| v(j) * <<i E I| <<k E K| 
            <<p E P| x(i,j,k,p) >> >> >>^2.5>>;
  lsum1(i E I, p E P): 
     <<j E J| <<k E K| x(i,j,k,p) >> >> =l= s(i,p);
  lsum2(j E J, k E K, p E P): 
     <<i E I| x(i,j,k,p) >> =E= d(k,p)*y(j,k);
  lsum3(j E J): <<k E K| 
     <<p E P| d(k,p)*y(j,k) >> >> =l= vu(j)*z(j);
  lsum4(j E J):  <<k E K| 
     <<p E P| d(k,p)*y(j,k) >> >> =g= vl(j)*z(j);
  lsum5(k E K): <<j E J| y(j,k)>> =e= 1;
}}
\end{verbatim}
\vfill
\end{slide}

\begin{slide}
\lhead{OPL 1/2}
\thispagestyle{fancy}
\tiny
\begin{verbatim}
// The network configuration
enum Sites ...;
enum Articles ...;

struct link {
    Sites    org;
    Sites    dst;
};

{link} Links = ...;

float+ cost[Links] = ...;
float+ siteCapacity[Sites] = ...;
float+ demand[Sites,Articles] = ...;
float+ linkCapacity[Links] = ...;
float+ onHand[Sites,Articles] = ...;

int minmax = 1;

var float+ totalCost;
var float+ siteExtra[Sites];
var float+ linkExtra[Links];
var float+ flow[Links,Articles];

minimize
   totalCost * minmax

\end{verbatim}
\vfill
\end{slide}

\begin{slide}
\lhead{OPL 2/2}
\thispagestyle{fancy}
\tiny
\begin{verbatim}
   
subject to {
   forall(l in Links)
      sum(a in Articles) 
         flow[l,a] <= linkCapacity[l] + linkExtra[l];

   forall(s in Sites)
      sum(a in Articles)   
         (sum(l in Links : s = l.dst) flow[l,a] -
             sum(l in Links : s = l.org) flow[l,a] +
             onHand[s,a] -
             demand[s,a] ) <= 
      siteCapacity[s] + siteExtra[s];

   totalCost = sum(l in Links & a in Articles) 
      cost[l] * flow[l,a];

   forall(s in Sites & a in Articles)
      demand[s,a] - onHand[s,a] <=   
         sum(l in Links : s = l.dst) flow[l,a] -
         sum(l in Links : s = l.org) flow[l,a];           
};

display linkExtra;
display siteExtra;
display flow;

\end{verbatim}
\vfill
\end{slide}

\begin{slide}
\lhead{Why \Zimpl}
\thispagestyle{fancy}
\normalsize
Why build yet another MPL?

\small
While there are over a dozend MPL's, all seem to be
commercial. None of them is AFAIK available with source code.

The are several possible reasons for this:

\begin{itemize}
\item MPL's are not a research topic\\
      (It's pretty clear what they possibly can do and what not)
\item There main purpose is to make modeling easier\\
      (Which increases the demand for model solving tools)
\item A lot MPL's started as research and went commercial\\
      (The same situation as with LP-Solvers)
\item They don't help if one is looking at a single model 
      that is at the edge of solvability\\
      (A common case in research)
\item By now all major suppliers of optimization software have a MPL
      bundled with there products\\
      (And let you download free \emph{student} versions)
\end{itemize}
Nonetheless having a platform independent 
non commercial open source MPL would have some advantages other than
the possibility to teach modeling with more than 300 variables.

\vfill
\end{slide}

\begin{slide}
\lhead{Model example}
\thispagestyle{fancy}
A model consists of 
%\begin{center}
\begin{tabular}{l}
\bf Sets\\
\bf Parameters\\
\bf Variables\\
\bf Objective\\
\bf Constraints\\
\end{tabular}

\vfill
\footnotesize
\begin{tabular}{lll}
$B$ & set of BSC's\\
$M$ & set of MSC's\\
$R$ & set of resources\\
$S$ & set of building stages\\
$A$ & $\subseteq B\cup M$ set of links between BSC's and MSC's\\\\
$c_a$ & $a\in A$ cost for connecting a BSC to a MSC\\
$f_s$ & $s\in S$ cost for building stage $s$\\
$d_{br}$ & $b\in B, r\in R$ demand of BSC $b$ for resource $r$\\
$k_{rs}$ & $r\in R, s\in S$ capacity of resource $r$ in stage $s$\\\\
$x_a$ & $\in\{0,1\}, a\in A$ 1 if link $a$ is used, 0 otherwise\\
$z_{ms}$ & $\in\{0,1\}, m\in M, s\in S$ 1 if MSC $m$ is in stage $s$\\
\end{tabular}
$$
\renewcommand{\arraystretch}{1.5}
\begin{array}{rlll}
\min&\displaystyle\sum_{bm\in A} c_{bm} x_{bm} + 
\sum_{m\in M}\sum_{s\in S} f_s z_{ms}\\
\mbox{s.t.}&\displaystyle\sum_{bm\in A} x_{bm}&=1&\displaystyle\forall b\in B\\
&\displaystyle\sum_{s\in S} z_{ms}&=1&\displaystyle\forall m\in M\\
&\displaystyle\sum_{bm\in A} d_{br} x_{bm} - \sum_{s\in S} k_{rs} z_{ms}
&\leq 0&\displaystyle\forall m\in M, r\in R
\end{array}
$$
\end{slide}

\begin{slide}
\lhead{\Zimpl example 1/2}
\thispagestyle{fancy}
{\tiny
\begin{verbatim}
set M := { "Altenburg", "Augsburg", "Berlin", ...
set B := { "Bischofsheim", "Bockenem", "Bonn", ...
set R := { "Lines", "Users" }; 
set S := { 0 to 6 };

set A := { 
   <"Berlin","Altenburg">, 
   <"Berlin","Braunschweig">,   
   <"Berlin","Chemnitz">, ...

param anbind_kosten[A] := 
   <"Aachen","Gie�en"> 479740, 
   <"Aachen","Hagen">  365308,
   <"Aachen","K�ln">   275908, ...

param bedarf[B * R] :=
 <"Bautzen", "Lines">      2, 
 <"Berlin",  "Lines">   1090, 
 <"Bautzen", "Users">    986,
 <"Berlin",  "Users"> 431252, ...

param kapazitaet[R * S] := 
 <"Lines", 0>     0, <"Users", 0>      0,
 <"Lines", 1> 18750, <"Users", 1> 200000,
 <"Lines", 2> 18750, <"Users", 2> 400000,
 <"Lines", 3> 37500, <"Users", 3> 600000,
 <"Lines", 4> 37500, <"Users", 4> 8000000,	
 <"Lines", 5> 75000, <"Users", 5> 1000000,	
 <"Lines", 6> 75000, <"Users", 6> 1200000;

param ausbau_kosten[S] := 
 <0>       0,  <1>  810000,  <2> 1782000,
 <3> 1924000,  <4> 2264000,  <5> 2614000,
 <6> 3002000;
\end{verbatim}
}
\vfill
\end{slide}
%------------------------------------------------------------------------------
\begin{slide}
\lhead{\Zimpl example 2/2}
\thispagestyle{fancy}
{\tiny
\begin{verbatim}
var x[A    ] binary;
var z[M * S] binary;

minimize kosten: 
   sum <b,m> in A     : anbind_kosten[b,m] * x[b,m] 
+  sum <m,s> in M * S : ausbau_kosten[s] * z[m,s];
  
# Jede BSC mu� an genau eine MSC angebunden werden.
#
subto c1: forall <b> in B do 
   sum <b,m> in A : x[b,m] == 1;

# Jede MSC mu� eine Konfiguration haben
#
subto c2: forall <m> in M do 
   sum <s> in S : z[m,s] == 1;

# Die Kapazit�t der gew�hlten Konfiguration 
# mu� ausreichen.
#
subto c3: forall <m,r> in M * R do 
   sum <b,m> in A : bedarf[b,r] * x[b,m] 
 - sum <s> in S : kapazitaet[r,s] * z[m,s] <= 0;

\end{verbatim}
}
\vfill
\end{slide}


\begin{slide}
\lhead{Examples}
\thispagestyle{fancy}
\tiny
\begin{verbatim}
set A := { 1, 2, 3 };
set B := { "hi", "ha", "ho" };
set C := { <1,2,"x">, <6,5,"y">, <787,12.6,"oh"> };
set D := A cross B;
set E := { 6 to 9 } union A without { <2>, <3> }; 
\end{verbatim}

\begin{verbatim}
param q := 5;
param u[A] := <1> 17, <2> 29, <3> 12;
param w[C] := <1,2,"x"> 1/2, <6,5,"y"> 2/3;
\end{verbatim}

\begin{verbatim}
var x1;
var x2 binary;
var y[A] real >= 2 <= 18;
var z[<a,b,c> in C] integer 
    >= a * 10 
    <= if b <= 3 then p[b] else 10 end; 
\end{verbatim}

\begin{verbatim}
minimize cost: 12 * x1 -4.4 * x2 
   + sum <a> in A : u[a] * y[a]
   + sum <a,b,c> in C with a in E and b > 3 : 
        -a/2 * z[a,b,c];
\end{verbatim}

\begin{verbatim}
subto time : 3 * x1 + 4 * x2 <= 7;
subto space: sum <a> in A : 2 * u[a] * y[a] >= 50;
subto weird: forall <a> in A : 
   sum <a,b,c> in C : z[a,b,c] == 55;
\end{verbatim}
\end{slide}

\begin{slide}
\lhead{\Zimpl}
\thispagestyle{fancy}

\normalsize What type of models can \Zimpl generate?

\small
LP, MIP

\normalsize Which solvers are supported?

\small
Since \Zimpl generates MPS or LPF files, it can be used with
nearly all available solvers. Integration within SIP is planned.

\normalsize How does it work?

\small
\Zimpl consist of two parts. The parser and the engine. The parser is
responsible for understanding the input and building an instruction
tree. The engine then traverses the tree and executes the instructions.

\normalsize Further directions

\small
\begin{itemize}
\item More possibilities for specifying data
\item Improved and extended syntax 
\item Hints to the solver
\item Commands to generate special inequalities 
\end{itemize}

\end{slide}


\end{document}

